Из обзора §2 видно, что среди обучаемых эмбеддингов \texttt{Persformer} наиболее выразителен, но его квадратичная по числу точек сложность плохо сочетается с многоканальным выходом PHT-фильтрации, где для каждого изображения формируется множество из сотен и тысяч точек. С другой стороны, \texttt{DeepSets} масштабируется линейно, но не моделирует попарные взаимодействия. Хочется получить архитектуру, которая (i) сохраняет выразительность self-attention, (ii) масштабируется приблизительно линейно по числу токенов диаграммы, (iii) умеет естественно использовать direction-метаинформацию точек, накопленную на этапе PHT. Эта мотивация в точности совпадает с задачей, для которой была предложена архитектура \texttt{Perceiver}~\cite{jaeglePerceiverGeneralPerception2021}: масштабируемая обработка больших унифицированных входов через cross-attention к компактному обучаемому латентному массиву. Мы адаптируем эту идею к задаче эмбеддинга диаграмм персистентности и называем получающуюся модель \texttt{LatentPersformer}.

\subsection{Концептуальное устройство}

\texttt{LatentPersformer} устроен как Perceiver-подобный энкодер с латентным bottleneck:

\begin{enumerate}
\item \textbf{Токенизация входа.} Каждая точка объединённой PHT-диаграммы изображения превращается в небольшой вектор-токен. В этот вектор входят: координаты $(b, d)$ точки, индикатор гомологической размерности, флаг типа фильтрации (sublevel или направленная PHT), направление $\alpha$ направленной фильтрации и его синусоидальное позиционное кодирование. Тем самым информация о направлении и о типе фильтрации зашита в признаки самих токенов, а не в позицию во входной последовательности. Линейный input embedding с LayerNorm и нелинейностью переводит токены в общее пространство модели.

\item \textbf{Cross-attention к латентам.} Вместо того чтобы выполнять self-attention напрямую над всеми $M$ токенами диаграммы (что стоило бы $\mathcal{O}(M^2)$), мы вводим небольшой обучаемый массив латентных векторов $L \in \mathbb{R}^{N \times d_L}$, $N \ll M$. Cross-attention использует латенты $L$ в качестве queries, а токены диаграммы --- в качестве keys и values. Эта операция распределяет содержимое произвольно большой диаграммы по фиксированному количеству <<слотов>> $N$ и имеет вычислительную сложность $\mathcal{O}(N \cdot M)$, то есть линейную по числу точек диаграммы.

\item \textbf{Self-attention над латентами.} Над массивом $L$ применяется стек обычных self-attention блоков (с residual connections, LayerNorm и MLP-слоями). Здесь self-attention уже дёшев: его сложность $\mathcal{O}(N^2)$ зависит только от выбранного размера $N$, не от длины входа. Глубина модели тем самым отвязывается от длины входа --- ключевое преимущество Perceiver-подобной декомпозиции.

\item \textbf{Pooling и предсказание.} После стека self-attention слоёв массив латентов агрегируется в один вектор. По умолчанию мы применяем \emph{attention-pool} с одной обучаемой query'ей $q \in \mathbb{R}^{d_L}$ (multi-head attention, см.~\cite{leeSetTransformer2019}): $\text{pooled} = \mathrm{MHA}(q, L, L)$. В отличие от среднего по латентам, attention-pool сохраняет стабильную идентичность отдельных слотов между разными входами, что критично для интерпретации латентов (см.~§\ref{ssec:exp-probes}). Альтернативная mean-pool ветка оставлена в реализации как ablation (см.~§\ref{ssec:exp-arch}). Полученный вектор подаётся в небольшой MLP-декодер, который выдаёт логиты классов (для классификации) или значение (для регрессии).
\end{enumerate}

\subsection{Свойства}

\paragraph{Перестановочная инвариантность.} В cross-attention не используется позиционное кодирование по индексу токена, а в самих токенах закодированы только содержательные признаки точки (включая направление $\alpha$). Поэтому перестановка точек диаграммы не меняет ни входное множество ключей/значений, ни результат softmax-агрегирования. Self-attention выполняется уже над фиксированным по структуре массивом латентов и не нарушает инвариантности.

\paragraph{Линейная сложность по числу точек.} В отличие от Persformer, у которого собственно self-attention имеет сложность $\mathcal{O}(M^2)$, в нашей архитектуре единственная зависимость от $M$ --- это cross-attention со сложностью $\mathcal{O}(N \cdot M)$. Это позволяет обрабатывать многоканальные PHT-диаграммы без выбрасывания части точек и без агрессивного субсэмплинга.

\paragraph{Декаплинг глубины и длины входа.} Все self-attention слои работают над фиксированным числом латентов $N$, поэтому глубина модели регулируется независимо от размера входа. На практике это означает, что мы можем строить модели заметно глубже, чем позволял бы Persformer на тех же диаграммах.

\paragraph{Учёт направления как часть фичи.} В Persformer и DeepSets направление PHT приходится подавать либо как дополнительный one-hot канал, либо моделировать отдельной диаграммой на каждое направление. В нашей токенизации направление $\alpha$ --- регулярный признак точки, и cross-attention сам выучивает, как соотносить токены из разных направлений через общий латентный массив. Это естественно отражает то обстоятельство, что разные направления PHT --- это \emph{разные взгляды на одну и ту же форму}.

\subsection{Сравнение с родственными моделями}

Связь с предыдущими работами можно резюмировать так:
\begin{itemize}
\item Относительно \texttt{Persformer}~\cite{reinauerPersformerTransformerArchitecture2021} мы заменяем самовнимание над диаграммой на cross-attention к фиксированному латенту, что обменивает выразительность одного слоя на масштабируемость и возможность глубоких стеков.
\item Относительно \texttt{DeepSets}~\cite{zaheerDeepSets2017} cross-attention реализует богаче, чем простую сумму: модель учится <<смотреть>> на разные подмножества точек диаграммы из разных латентных слотов.
\item Относительно \texttt{Perceiver}~\cite{jaeglePerceiverGeneralPerception2021} специфика в том, что роль <<byte array>> играет тонко токенизированное мультимножество точек диаграммы с направленной фильтрации, а не пиксели или сэмплы аудио, а позиционное кодирование носит не координатный (как в изображениях), а угловой (по направлению PHT) характер.
\end{itemize}

\begin{figure}[ht]
\centering
\fbox{\parbox{0.9\linewidth}{\centering [PLACEHOLDER]\\[0.3em]
Схема архитектуры \texttt{LatentPersformer}. Слева снизу: PHT-диаграммы изображения, объединённые в одно мультимножество точек, превращаются в последовательность токенов (каждый токен --- небольшой вектор-признак точки с направлением и типом фильтрации). Стрелка вверх ведёт в блок cross-attention: queries --- обучаемый латентный массив $L \in \mathbb{R}^{N \times d_L}$, keys/values --- токены диаграммы. Над cross-attention --- стек self-attention слоёв над $L$. Завершают пайплайн attention-pool с одной обучаемой query $q$ (по умолчанию) или усреднение по латентам (mean-pool, опция) и MLP-голова, выдающая предсказание.}}
\caption{Архитектура \texttt{LatentPersformer}.}
\label{fig:latent-persformer}
\end{figure}
